\documentclass[tikz,border=10pt]{standalone}
\usepackage{ctex}
\usepackage{tikz}
\usetikzlibrary{arrows.meta,positioning}

\begin{document}
\begin{tikzpicture}[
  font=\small,
  >=Latex,
  bar/.style={draw=black!70, rounded corners=1.5mm, minimum height=8mm, align=center},
  ax/.style={very thick, -Latex}
]

% axis
\draw[ax] (0,0) -- (12.2,0) node[right] {相对开销 / 时延};
\draw[ax] (0,0) -- (0,4.2) node[above] {操作类型};

% bars
\node[anchor=east] at (-0.2,3.2) {写入 \texttt{push}};
\node[bar, fill=green!25, minimum width=2.2cm, anchor=west] at (0.2,3.2) {\(O(1)\)};

\node[anchor=east] at (-0.2,2.0) {告警触发检索};
\node[bar, fill=orange!25, minimum width=5.4cm, anchor=west] at (0.2,2.0) {最坏 \(O(C)\)};

\node[anchor=east] at (-0.2,0.8) {常规无告警周期};
\node[bar, fill=blue!25, minimum width=2.8cm, anchor=west] at (0.2,0.8) {仅写入路径};

\node[draw=black!35, rounded corners=2mm, fill=teal!10, align=left, inner sep=5pt] at (8.3,3.3) {
实时性结论：\\
1) 主路径由 \(O(1)\) 写入主导；\\
2) 检索开销按告警触发；\\
3) 固定容量保证内存上限可控。
};

\end{tikzpicture}
\end{document}
